\documentclass[10pt]{article}
\usepackage[T1]{fontenc}
\usepackage[utf8]{inputenc}
\usepackage[french]{babel}
\usepackage{amsmath,amssymb,mathrsfs,mathpazo}
\usepackage{geometry}
\geometry{a4paper,margin=2cm}
\usepackage{graphicx,array,tabularx,tabularray,colortbl,multirow,multicol,booktabs,enumitem,pifont,ulem,fancybox,diagbox,cancel}
\usepackage{pstricks,pstricks-add,pst-plot,pst-tree,pst-node,pst-text,pst-func,pst-math,pst-bspline,pst-3dplot,pst-eucl,pst-solides3d}
\let\pscalculate\undefined
\usepackage{pst-bezier}
\usepackage{tikz,pgfplots,tkz-tab}
\usetikzlibrary{arrows,arrows.meta,patterns,calc,positioning,trees,decorations.markings,shapes.geometric,shapes.misc,shapes.arrows,intersections,decorations.pathmorphing,decorations.pathreplacing,backgrounds,shadows,babel}
\pgfplotsset{compat=1.18}
\usepgfplotslibrary{fillbetween}
\usepackage[np]{numprint}
\usepackage{listings,makecell,pas-tableur,twoopt}
\let\barre\undefined
\usepackage[pstricks]{tablvar}
\definecolor{darkgreen}{RGB}{0,100,0}
\providecommand{\red}{\color{red}}
\providecommand{\blue}{\color{blue}}
\providecommand{\green}{\color{green}}
\providecommand{\black}{\color{black}}
\providecommand{\white}{\color{white}}
\providecommand{\gray}{\color{gray}}
\providecommand{\cyan}{\color{cyan}}
\providecommand{\magenta}{\color{magenta}}
\providecommand{\yellow}{\color{yellow}}
\providecommand{\violet}{\color{violet}}
\providecommand{\orange}{\color{orange}}
\providecommand{\euro}{\texteuro}
\providecommand{\EUR}{\texteuro}
\providecommand{\ieme}{\textsuperscript{e}}
\providecommand{\degres}{\ensuremath{^\circ}}
\providecommand{\degree}{\ensuremath{^\circ}}
\providecommand{\vv}[1]{\overrightarrow{#1}}
\providecommand{\Coord}[1]{#1}
\providecommand{\decosix}{}
\providecommand{\starredbullet}{\ensuremath{\bullet}}
\pagestyle{empty}
\setlength{\parindent}{0pt}
\usepackage{adjustbox}\DeclareUnicodeCharacter{2236}{:}\DeclareUnicodeCharacter{2190}{\ensuremath{\leftarrow}}
\providecommand{\R}{\mathbb{R}}
\providecommand{\N}{\mathbb{N}}
\providecommand{\Z}{\mathbb{Z}}
\providecommand{\Q}{\mathbb{Q}}
\providecommand{\D}{\mathbb{D}}
\providecommand{\C}{\mathbb{C}}
\providecommand{\e}{{\mathrm{e}}}
\providecommand{\vect}[1]{\overrightarrow{#1}}
\providecommand{\vectt}[1]{\overrightarrow{\mathrm{#1}}}
\providecommand{\vvt}[1]{\overrightarrow{\mathrm{#1}}}
\providecommand{\np}[2][]{{#2~#1}}
\providecommand{\vv}[1]{{\overrightarrow{#1}}}
\renewcommand{\d}{{\mathrm{d}}}
\providecommand{\degres}{{^\circ}}
\providecommand{\degre}{{^\circ}}
\providecommand{\blue}{\color{blue}}
\providecommand{\red}{\color{red}}
\providecommand{\enskip}{\quad}
\providecommand{\llbracket}{\lbrack\!\lbrack}
\providecommand{\rrbracket}{\rbrack\!\rbrack}
\providecommand{\euro}{{\text{€}}}
\providecommand{\no}{{\text{n°}}}
\providecommand{\interval}[2]{{[#1;#2]}}
\providecommand{\intervFO}[2]{{[#1;#2[}}
\providecommand{\dotfill}{\ldots}

\providecommand{\ds}{\displaystyle}
\providecommand{\textding}[1]{#1}
\providecommand{\Oij}{\(\mathrm{O};\vec{\imath},\vec{\jmath}\)}
\providecommand{\Oijk}{\(\mathrm{O};\vec{\imath},\vec{\jmath},\vec{k}\)}
\providecommand{\Ouv}{\(\mathrm{O};\vec{u},\vec{v}\)}
\providecommand{\intervFF}[2]{\ensuremath{[#1;#2]}}
\renewcommand{\euro}{\ifmmode\text{\texteuro}\else\texteuro\fi}
\DeclareUnicodeCharacter{00B0}{\ifmmode^{\circ}\else\textdegree\fi}
\DeclareUnicodeCharacter{20AC}{\ifmmode\text{\texteuro}\else\texteuro\fi}
% Énoncés : Éducation nationale. Transcription et corrigés : APMEP, Denis Vergès et contributeurs. Découpage et classement : CoopMaths, Sébastien Lozano, Éric Elter et contributeurs. Adaptation de la présentation : Maths Explorer.
% CC BY-SA : https://coopmaths.fr/www/about/
% https://coopmaths.fr/alea/static/eam/2026/tex/eam_specifique_2026_06_metropole_2.tex
\begin{document}
\noindent\textit{APMEP / CoopMaths. CC BY-SA. Présentation adaptée par Maths Explorer.}\par\medskip
En 2025, Emma et Pierre disposent d'un capital initial de $\np{20000} $\,\euro{}.

	Pour réaliser un projet, ils ont besoin que le capital s'élève à $\np{22000} $\,\euro{}.

	Ils souhaitent donc placer leur capital durant plusieurs années.

	Ils ont le choix entre deux placements :

	\subsection*{A -- Premier placement.}
	Avec ce placement A, le capital augmente chaque année de $200 $\,\euro{} par rapport au capital de l'année précédente.

	On note $a_{n}$ la somme disponible en $2025+n$.

	On a donc $a_{0}=\np{20000}$.

	\begin{enumerate}
		\item Calculer $a_{1}$ et $a_{2}$\,.
		\item \begin{enumerate}
			\item Exprimer $a_{n+1}$ en fonction de $a_{n}$ pour tout entier naturel $n$.

			\item En déduire la nature de la suite $\left(a_{n}\right)$.

			Préciser sa raison.
		\end{enumerate}

		\item Exprimer $a_{n}$ en fonction de $n$ pour tout entier naturel $n$.

		\item À partir de quelle année auront-ils la somme nécessaire pour leur projet ?
	\end{enumerate}

	\subsection*{B -- Second placement.}
	Avec le placement B, le capital augmente chaque année de $2\,\%$ par rapport au capital de l'année précédente.

	On note $b_{n}$ la somme disponible en $2025+n$.

	On a donc $b_{0}=\np{20000}$.

	\begin{enumerate}
		\item Détailler le calcul qui permet d'obtenir $b_{1}=\np{20400}$.

		\item \begin{enumerate}
			\item Exprimer $b_{n+1}$ en fonction de $b_{n}$ pour tout entier naturel $n$.

			\item En déduire la nature de la suite $\left(b_{n}\right)$.

			Préciser sa raison.
		\end{enumerate}
		\item Exprimer $b_{n}$ en fonction de $n$ pour tout entier naturel $n$.
		\item \begin{minipage}[t]{0.6\linewidth}
			Grâce à l'extrait de tableur ci-contre, on visualise le capital arrondi à l'euro près.

		Déterminer à partir de quelle année ils auront la somme nécessaire pour leur projet.
		\end{minipage}
		\hfill
		\begin{tblr}[T]{colspec = {X[1cm,c]X[2cm,c]},hlines, vlines,cell{2-Z}{2}={cmd=\np}}
			$n$ & $b_{n}$ \\
			0   & 20000   \\
			1   & 20400   \\
			2   & 20808   \\
			3   & 21224   \\
			4   & 21649   \\
			5   & 22082   \\
			6   & 22523   \\
			7   & 22974   \\
			8   & 23433   \\
			9   & 23902   \\
			10  & 24380
		\end{tblr}
	\end{enumerate}

\subsection*{C -- Bilan}
	Quel placement conseiller à Emma et Pierre afin qu'ils puissent réaliser leur projet le plus tôt possible ? Justifier.

\end{document}
